\chapter{Theory and Problem Formulation}
\label{chap:theory}

This chapter presents the concepts needed to formulate CT-based mucus-plug prediction as a machine-learning problem. A chest CT examination is treated as volumetric image data composed of multiple local observations, while the target is defined at patient level. This creates a weakly supervised regression setting in which local image evidence must be represented, modelled, and related to one scan-level output.


\section{CT-Based Prediction Problem}
\label{sec:ct_based_prediction_problem}

\subsection{Volumetric CT as Structured Input Data}
\label{sec:volumetric_ct_input}

CT-based prediction problems differ from standard two-dimensional image prediction tasks because the input is a volumetric medical examination rather than a single isolated image. A chest CT scan consists of a sequence of cross-sectional slices that together represent thoracic anatomy \cite{whiting2015ctbasic,dejong2005ctairways}. For a patient or examination indexed by \(i\), the CT volume is represented as

\[
X_i = \{x_{i,1}, x_{i,2}, \dots, x_{i,n_i}\},
\]

where \(X_i\) denotes the full CT examination, \(x_{i,j}\) denotes the \(j\)-th axial slice or local image instance, and \(n_i\) is the number of slices in the volume. Relevant information may be distributed across several slices, localized to a small part of the volume, or absent from large regions of the scan \cite{xue2023copdmil,djahnine2024weakly}. The scan must therefore be treated as a structured set of local observations rather than as one flat image.

\subsection{Patient-Level Targets and Regression Formulation}
\label{sec:patient_level_regression}

In CT-based mucus-plug assessment, the clinically meaningful quantity is commonly summarized as an overall burden at the patient or scan level, rather than as an independent label for each individual slice \cite{dunican2018mucus,dunican2018autopsy}. Let

\[
y_i \in \mathbb{R}
\]

denote the patient-level target associated with examination \(X_i\). The prediction task can then be written as

\[
\hat{y}_i = f(X_i),
\]

where \(f(\cdot)\) is a prediction function that maps the CT examination to an estimated patient-level output \(\hat{y}_i\). Since the target represents a burden measure rather than a discrete class label, the problem is naturally formulated as regression:

\[
f : X_i \mapsto \hat{y}_i \in \mathbb{R}.
\]

This formulation defines both supervision and evaluation at the scan level. Individual slices may contain useful local evidence, but the final prediction is meaningful only as a patient-level quantity.


\subsection{Weak Supervision and Label Granularity}
\label{sec:weak_supervision_label_granularity}
A central theoretical difficulty is the mismatch between the level at which image evidence is observed and the level at which supervision is available. The training data can be written as

\[
\mathcal{D} = \{(X_i, y_i)\}_{i=1}^{N},
\]

where each CT examination \(X_i\) is associated with one patient-level target \(y_i\). However, corresponding slice-level labels

\[
y_{i,j}
\]

are not observed for the individual slices \(x_{i,j}\). The model therefore has access to the overall target for the examination, but not to annotations indicating which slices, regions, or anatomical locations are responsible for that target. This makes the problem weakly supervised, the labels are coarser than the image evidence from which the model must learn \cite{ren2023weakly}. Related CT studies address similar situations by treating a scan as a collection of local instances and combining instance-level evidence into a case-level prediction \cite{qi2021drmil,xue2023copdmil}. The same issue arises here, although the target is a continuous burden score rather than a class label.


\subsection{Implications for CT-Based Learning}
\label{sec:implications_ct_learning}

The task can therefore be understood as patient-level regression under weak supervision. It requires a suitable CT representation, feature extraction from local image evidence, aggregation from slice-level information to a scan-level output, and evaluation at patient level. Limited cohort size makes these choices more important, because highly flexible models may not be supported by the available supervision \cite{dhiman2023samplesize,litjens2017survey}. For volumetric medical imaging, the choice between slice-based, instance-based, and fully volumetric formulations is shaped by annotation level, computational cost, transfer-learning possibilities, and generalization risk \cite{wang2023twoseventyfive}.








\section{Preprocessing of CT Data}
\label{sec:preprocessing_ct_data}

\subsection{Purpose of Preprocessing in CT-Based Learning}
\label{sec:purpose_preprocessing_ct}

Preprocessing transforms CT data into a form suitable for computational modelling. In CT-based learning, this is needed because medical images may vary in scanner settings, acquisition protocols, reconstruction choices, and anatomical coverage \cite{rizzo2018radiomics,litjens2017survey}.

At a general level, preprocessing can be written as a transformation
\[
\tilde{X}_i = T(X_i),
\]
where \(X_i\) denotes the original CT examination, \(T(\cdot)\) denotes a preprocessing operation or sequence of operations, and \(\tilde{X}_i\) denotes the transformed examination used in later analysis. When preprocessing is considered at slice level, the same idea can be written as
\[
\tilde{x}_{i,j} = T(x_{i,j}),
\]
where \(x_{i,j}\) is an axial CT slice and \(\tilde{x}_{i,j}\) is the corresponding preprocessed slice.

The purpose of preprocessing is not to solve the prediction problem directly.
Rather, preprocessing defines how the raw image evidence is made available for later analysis. It may reduce technical variation, impose numerical consistency, or emphasize image structures that are relevant for subsequent representation learning. In digital image processing, such operations are commonly used to improve the suitability of image data for later analysis, for example by adjusting intensity ranges, reducing unwanted variation, or emphasizing selected spatial structures \cite{gonzalez2018digital}.


\subsection{Intensity Transformation and CT Windowing}
\label{sec:intensity_windowing_ct}

CT images are represented in terms of attenuation values, commonly expressed in Hounsfield units \cite{whiting2015ctbasic}. These values cover a wide numerical range, so a restricted intensity mapping is often needed before visual interpretation or computational analysis.

Windowing is a standard way of mapping CT attenuation values into a more focused intensity range. In clinical CT interpretation, different windows are used to emphasize different anatomical or tissue relationships, such as air-filled lung structures or denser soft-tissue regions \cite{whiting2015ctbasic}. More generally, windowing can be understood as an intensity transformation applied before later analysis. For a slice \(x_{i,j}\), a windowed representation can be written as
\[
x_{i,j}^{(w)} = T_w(x_{i,j}),
\]
where \(T_w(\cdot)\) denotes the intensity transformation associated with window setting \(w\).

Windowing does not change the anatomy in the slice; it changes which attenuation relationships are emphasized. This matters in chest CT because lung parenchyma, airways, vessels, and soft-tissue structures are not equally visible under one intensity mapping \cite{dejong2005ctairways}. From a machine-learning perspective, CT windowing is therefore part of the input representation problem. Prior work on multi-window CT learning supports the general idea that different CT windows can provide complementary information, although the choice remains task dependent \cite{lu2020multiwindow}.


\subsection{Spatial Standardization and Slice-Level Inputs}
\label{sec:spatial_standardization_slice_inputs}

CT examinations can also differ in matrix size, field of view, slice thickness, slice spacing, number of slices, and anatomical coverage \cite{rizzo2018radiomics}. These differences create challenges for models that require inputs with consistent numerical dimensions and comparable spatial organization \cite{litjens2017survey}.

Spatial preprocessing addresses this issue by transforming the image data into a more standardized form. Depending on the modelling formulation, this may involve operations such as resizing, resampling, cropping, padding, or restricting the analysis to a defined image region. Abstractly, a spatial preprocessing operation can be included in the same transformation notation
\[
\tilde{x}_{i,j} = T_{\mathrm{spatial}}(x_{i,j}),
\]
where \(T_{\mathrm{spatial}}(\cdot)\) denotes a transformation that changes the spatial format of the slice while aiming to preserve relevant anatomical content.

Different modelling choices require different spatial preparation. A volumetric model, a slice-based model, and a patch-based model may not use the same input units, but each requires image data that can be processed consistently by the learning method \cite{wang2023twoseventyfive}. In slice-level formulations, axial slices act as local observations within the larger examination. This does not make the final task slice-level prediction; supervision remains defined for the full examination \cite{ren2023weakly,xue2023copdmil}.

\subsection{Candidate Slice Relevance}
\label{sec:candidate_slice_relevance}

A volumetric CT examination may contain many slices, but not all slices are equally informative for a scan-level prediction task. This is particularly important in weakly supervised settings, where the label describes the case as a whole but does not identify which local observations are responsible for that label \cite{ren2023weakly}. At a general level, one may describe a candidate set of local observations as
\[
C_i \subseteq X_i,
\]
where \(C_i\) denotes the set of slices or image instances retained as candidate evidence from the full CT examination \(X_i\). This notation is deliberately general: it does not specify how the candidate set is obtained, only that a CT-based learning problem may involve deciding which parts of the volume are considered suitable for later analysis.

Candidate selection can reflect anatomical relevance, image suitability, and redundancy reduction. Related weakly supervised and multiple-instance CT studies highlight that local evidence in volumetric images is often heterogeneous rather than uniformly informative \cite{xue2023copdmil,djahnine2024weakly}. The exact method used to filter or select slices is methodological; the theoretical point is that a volumetric scan contains many local observations whose relevance to the patient-level target may vary.



\section{Feature Extraction from CT Images}
\label{sec:feature_extraction_ct_images}

\subsection{Purpose of Feature Extraction}
\label{sec:purpose_feature_extraction}

Feature extraction transforms image data into representations that make relevant visual information available for prediction. In CT-based learning, the raw image array is high-dimensional, so a feature extractor maps a preprocessed image or region into a more compact feature space.

For a preprocessed CT slice \(\tilde{x}_{i,j}\), feature extraction can be written generally as
\[
z_{i,j} = \phi(\tilde{x}_{i,j}),
\]
where \(\phi(\cdot)\) denotes a feature-extraction function and \(z_{i,j}\) is the resulting feature representation. The notation is deliberately general. The function \(\phi(\cdot)\) may represent a handcrafted descriptor, a learned neural-network mapping, or another transformation that converts image content into a representation suitable for modelling.

Preprocessing prepares the image data, while feature extraction defines which visual properties are made available for prediction. In scan-level CT prediction, useful features should preserve local evidence that may contribute to the patient-level target while reducing variation that is unlikely to help prediction \cite{rizzo2018radiomics,litjens2017survey}.



\subsection{Handcrafted Image Features}
\label{sec:handcrafted_image_features}

Handcrafted image features are predefined descriptors that encode visual properties such as intensity, texture, edges, gradients, local patterns, or spatial relationships. In medical imaging, such features have often been used to describe image structure compactly, especially when labelled data are limited \cite{castellano2004texture,rizzo2018radiomics}.

A general handcrafted feature extractor can be written as
\[
z_{i,j}^{(\mathrm{hand})} = \phi_{\mathrm{hand}}(\tilde{x}_{i,j}),
\]
where \(\phi_{\mathrm{hand}}(\cdot)\) denotes a predefined feature function. The resulting vector \(z_{i,j}^{(\mathrm{hand})}\) may encode local texture, intensity statistics, edge information, or other manually specified image characteristics. Such features are not optimized end-to-end for the target task, but they may provide useful summaries of visual structure when the relevant image patterns are expected to be local or when the available training data are limited.

Their strength is that they encode prior assumptions about useful image properties, which can be valuable with limited labels. Their limitation is that the representation is fixed before learning; if the descriptors miss the relevant evidence, the downstream model cannot recover it.



\subsection{Learned Feature Representations}
\label{sec:learned_feature_representations}

Learned feature representations are estimated from data rather than defined manually. In deep learning, feature extraction becomes part of the learning process rather than a separate fixed descriptor stage \cite{lecun1998gradient,litjens2017survey}.

A learned feature extractor can be written as
\[
z_{i,j}^{(\mathrm{learned})} = \phi_{\theta}(\tilde{x}_{i,j}),
\]
where \(\phi_{\theta}(\cdot)\) is a parameterized mapping with trainable parameters \(\theta\). Unlike handcrafted descriptors, the representation is shaped by the data, the objective function, and the available supervision.

For CT images, learned representations are useful because relevant evidence may be complex, spatially distributed, and difficult to specify manually. Their flexibility is also their main risk: it must be supported by enough data and appropriate supervision, as discussed in Section~\ref{sec:model_capacity_limited_data}.




\subsection{Feature Dimensionality and Representation Quality}
\label{sec:feature_dimensionality_representation_quality}

The quality of a feature representation depends on both its information content and its dimensionality. A useful representation should preserve target-relevant image properties while reducing redundant or unrelated variation. This balance is especially important in medical imaging, where image data are high-dimensional and patient cohorts are often limited \cite{mwangi2014feature,rizzo2018radiomics}.

If a representation is too low-dimensional, relevant image information may be lost. If it is too high-dimensional relative to the number of labelled cases, the downstream model may become more vulnerable to overfitting and unstable estimation. Dimensionality reduction and feature selection can therefore be understood as general strategies for improving the match between the information content of the representation and the amount of supervision available for learning \cite{mwangi2014feature}.

Handcrafted and learned representations encode image information according to different assumptions and may capture complementary aspects of medical images \cite{jeong2024handcrafteddeep,lai2018featurefusion}. Whether this improves prediction is empirical, but conceptually it shows why feature extraction determines what kind of evidence later model stages can use.



\section{CNN Modelling and Transfer Learning}
\label{sec:cnn_modelling_transfer_learning}

This section introduces CNNs, transfer learning, and fine-tuning as concepts needed to understand the implemented CT prediction pipeline under limited supervision.


\subsection{Convolutional Neural Networks for CT Images}
\label{sec:cnn_ct_images}

Convolutional neural networks are neural models designed for structured image data. They apply learnable filters across local image regions, allowing similar visual patterns to be detected across spatial locations \cite{lecun1998gradient,litjens2017survey}.

For a preprocessed CT slice \(\tilde{x}_{i,j}\), a CNN-based mapping can be written abstractly as
\[
z_{i,j} = \phi_{\theta}(\tilde{x}_{i,j}),
\]
where \(\phi_{\theta}(\cdot)\) denotes a convolutional feature extractor with trainable parameters \(\theta\), and \(z_{i,j}\) denotes the resulting learned representation. A task-specific prediction function may then map this representation to an instance-level output,
\[
s_{i,j} = h_{\theta}(z_{i,j}).
\]
This notation separates representation learning from prediction without committing to a specific architecture.

CNNs are relevant for CT images because useful evidence may be local, spatially structured, and difficult to summarize using global intensity statistics alone. However, CNN modelling does not remove the weakly supervised nature of the task: local representations still have to be related to patient-level targets through aggregation and evaluation.


\subsection{Residual Learning}
\label{sec:residual_learning}

Residual learning allows a network block to learn a residual transformation relative to its input rather than a completely new representation from scratch \cite{he2016residual}. This is relevant because deeper CNNs can represent richer features but may also be harder to optimize.

A residual block can be written abstractly as
\[
u_{l+1} = u_l + F_l(u_l; \theta_l),
\]
where \(u_l\) is the input representation at layer or block \(l\), \(F_l(\cdot)\) is a learnable residual function, \(\theta_l\) denotes the parameters of that block, and \(u_{l+1}\) is the resulting output representation. The additive identity connection allows information from earlier layers to pass forward directly, while the residual function learns how the representation should be modified.

For this thesis, residual learning matters mainly as the architectural principle behind the ResNet backbone used later. The choice of a concrete backbone remains a methodological decision.


\subsection{Model Capacity in Limited-Data Settings}
\label{sec:model_capacity_limited_data}

Model capacity refers to the flexibility of a model to represent complex functions. In CNNs, capacity is influenced by network depth, number of parameters, feature dimensionality, and prediction-head complexity. Higher-capacity models can represent more complex relationships, but require sufficient supervision to estimate their parameters reliably.

This issue is particularly important in medical imaging, where labelled datasets are often smaller than in general computer vision and where labels may be expensive or difficult to obtain \cite{litjens2017survey}. At a general level, model learning can be expressed as
\[
\theta^{*} = \arg\min_{\theta} \mathcal{L}_{\mathrm{train}}(\theta).
\]
When the model is highly flexible and the dataset is small, minimizing the training loss alone may not ensure good generalization.

Insufficient sample size can lead to unstable estimates and overly optimistic performance assessment \cite{dhiman2023samplesize}. In volumetric medical imaging, this concern is amplified by high-dimensional input and by the choice between two-dimensional, slice-based, and fully volumetric formulations \cite{wang2023twoseventyfive}. Model capacity is therefore part of the theoretical constraint imposed by limited medical image data.


\subsection{Transfer Learning as Representation Reuse}
\label{sec:transfer_learning_representation_reuse}

Transfer learning addresses limited-data settings by reusing information learned from a source task or source dataset. Instead of initializing all model parameters randomly, a model can begin from parameters already learned elsewhere \cite{kim2022transfer}.

In general notation, pretraining can be described as learning parameters on a source dataset,
\[
\theta_{\mathrm{source}}^{*}
=
\arg\min_{\theta}
\mathcal{L}_{\mathrm{source}}(\theta),
\]
and then using those parameters to initialize learning on the target task,
\[
\theta_0 \leftarrow \theta_{\mathrm{source}}^{*}.
\]
Transfer learning can therefore be understood as representation reuse: the pretrained model provides an initial feature space that can be adapted to the target prediction problem.

Transfer learning is relevant in medical imaging because target datasets are often limited, while deep models may require many labelled examples \cite{litjens2017survey,kim2022transfer}. Pretraining may use natural-image data or medical-image data; RadImageNet is one radiology-domain example \cite{mei2022radimagenet}. Its usefulness depends on the relationship between the source and target domains, so it should be understood as a strategy for limited supervision rather than a guaranteed improvement.


\subsection{Fine-Tuning and Domain Adaptation}
\label{sec:fine_tuning_domain_adaptation}

Fine-tuning is the process of adapting a pretrained model to a target task. After initialization from source-domain parameters, the model is trained further using target-domain data. This can be written abstractly as
\[
\theta_{\mathrm{target}}^{*}
=
\arg\min_{\theta}
\mathcal{L}_{\mathrm{target}}(\theta; \theta_0),
\]
where \(\theta_0\) denotes the pretrained initialization and \(\mathcal{L}_{\mathrm{target}}\) is the target-task objective.

Fine-tuning strategies differ in how much of the pretrained representation is reused. Keeping most layers fixed treats the model mainly as a feature extractor, while updating more layers allows stronger adaptation to the target domain. This balances stability and flexibility: preserving pretrained features may reduce overfitting when target data are limited, while adapting more parameters may be needed when source and target domains differ \cite{kim2022transfer}. The exact fine-tuning procedure is methodological; the theoretical role is to align a reused representation with the target prediction problem.



\section{CT Slice Aggregation}
\label{sec:ct_slice_aggregation}

\subsection{From Local Slice Evidence to Scan-Level Output}
\label{sec:local_slice_to_scan_output}

In scan-level CT prediction, the target is defined for the full examination, while much of the image evidence is observed locally at slice or region level. Aggregation is therefore required to combine local evidence into a single patient-level prediction. Related CT learning work uses similar formulations by treating an examination as a collection of slice-level instances that support a case-level output \cite{qi2021drmil,xue2023copdmil}.

Let \(s_{i,j}\) denote an instance-level score or evidence value derived from slice \(x_{i,j}\). A general aggregation function can be written as
\[
\hat{y}_i = A(s_{i,1}, s_{i,2}, \dots, s_{i,n_i}),
\]
where \(A(\cdot)\) maps local evidence values to a scan-level prediction \(\hat{y}_i\). If only a candidate subset of local observations is considered, the same idea can be written as
\[
\hat{y}_i = A\bigl(\{s_{i,j} : x_{i,j} \in C_i\}\bigr),
\]
where \(C_i \subseteq X_i\) denotes the local observations considered for prediction.

Different aggregation functions express different assumptions about how relevant evidence is distributed across the scan. In volumetric CT, disease-related evidence may be spatially heterogeneous rather than uniformly present throughout the volume \cite{xue2023copdmil}. Aggregation is therefore a conceptual mechanism for relating local CT evidence to a scan-level target.


\subsection{Multiple-Instance Learning Perspective}
\label{sec:mil_perspective}

The aggregation problem can be interpreted through multiple-instance learning. In this view, an input case is a bag of instances, while supervision is available only at bag level. For CT, the scan is the bag and slices or local regions are instances. The available label \(y_i\) describes the bag as a whole, while instance-level labels are unobserved.

Using this terminology, the CT examination
\[
X_i = \{x_{i,1}, x_{i,2}, \dots, x_{i,n_i}\}
\]
is a bag of local observations. Multiple-instance CT studies use this type of formulation to relate local image information to case-level outputs \cite{qi2021drmil,xue2023copdmil}. Although many multiple-instance studies address classification, the same conceptual structure applies to regression: local evidence must be combined into a continuous scan-level estimate rather than a discrete class probability.

This perspective separates two questions: feature extraction concerns how local evidence is represented, while aggregation concerns how that evidence is combined into a scan-level output.


\subsection{Aggregation Assumptions}
\label{sec:aggregation_assumptions}

Aggregation functions encode assumptions about how relevant evidence is distributed across the scan. A diffuse-evidence assumption treats many local observations as informative, while a focal-evidence assumption treats the target as driven mainly by a smaller number of highly informative observations \cite{qi2021drmil,xue2023copdmil}. Mixed or weighted assumptions allow both broad and localized evidence to influence the final prediction. The exact aggregation rule is methodological, but scan-level prediction always requires some assumption about how local CT evidence relates to the patient-level target.


\subsection{Aggregation under Weak Supervision}
\label{sec:aggregation_weak_supervision}

Aggregation becomes especially important under weak supervision because the model does not receive direct feedback about which slices are relevant. An accurate scan-level prediction does not necessarily mean that the model has identified the clinically relevant slices or regions. Different local evidence patterns may lead to similar scan-level outputs, so aggregation supports prediction but does not by itself solve localization or explanation \cite{djahnine2024weakly,almakky2025weakly}. Specific aggregation functions, hyperparameters, and empirical comparisons belong to the methodology and experiments.



\section{Regression and Evaluation}
\label{sec:regression_evaluation}

\subsection{Continuous Target Prediction}
\label{sec:continuous_target_prediction}

The prediction task is formulated as regression because the target is a numerical burden measure rather than a discrete class label. For a CT examination \(X_i\), a parameterized model produces a patient-level prediction
\[
\hat{y}_i = f_{\theta}(X_i),
\]
where \(y_i \in \mathbb{R}\) denotes the corresponding target. The prediction error for patient \(i\) can be written as
\[
e_i = y_i - \hat{y}_i.
\]

This formulation preserves the ordering and magnitude of the target. Errors are interpreted by their numerical size rather than only as correct or incorrect, which is appropriate when burden magnitude is meaningful.


\subsection{Regression Loss Functions}
\label{sec:regression_loss_functions}

Regression models are trained by minimizing a loss function that penalizes disagreement between observed and predicted target values:
\[
\theta^{*}
=
\arg\min_{\theta}
\mathcal{L}(\theta).
\]
A common objective for continuous prediction is mean squared error,
\[
\mathcal{L}_{\mathrm{MSE}}(\theta)
=
\frac{1}{N}
\sum_{i=1}^{N}
\left(y_i - \hat{y}_i\right)^2.
\]

The squared-error form penalizes larger deviations more strongly than smaller deviations. Other regression losses and metrics emphasize different aspects of error, so the optimization loss and the reported evaluation metrics should be distinguished \cite{terven2025lossmetrics,kocak2025medicalmetrics}.


\subsection{Patient-Level Evaluation}
\label{sec:patient_level_evaluation}

Evaluation must be performed at the same level as the target. In scan-level CT prediction, performance is assessed by comparing patient-level predictions with patient-level targets,
\[
\{(y_i, \hat{y}_i)\}_{i=1}^{N}.
\]
Intermediate slice-level values may be useful inside the model, but they should not be treated as independently validated outputs unless corresponding slice-level labels are available.

Error-based metrics summarize the magnitude of prediction errors, while association-based metrics describe how predictions vary with observed targets. Correlation coefficients measure association between continuous variables, but they do not by themselves describe absolute error or calibration \cite{schober2018correlation}. Patient-level regression evaluation should therefore use metrics that reflect the intended meaning of prediction error.


\subsection{Validation and Generalization in Small Cohorts}
\label{sec:validation_generalization_small_cohorts}

Validation estimates how well a model is expected to perform on unseen cases. Performance should be evaluated on data not used to fit the model parameters. If \(\mathcal{D}_{\mathrm{train}}\) denotes the training data and \(\mathcal{D}_{\mathrm{val}}\) denotes held-out validation data, learning and validation can be separated as
\[
\theta^{*}
=
\arg\min_{\theta}
\mathcal{L}_{\mathrm{train}}(\theta; \mathcal{D}_{\mathrm{train}}),
\]
with performance estimated on \(\mathcal{D}_{\mathrm{val}}\).

This separation is especially important in small-cohort medical imaging, where high-dimensional inputs and flexible models can produce overly optimistic performance estimates if evaluation is not carefully separated from training \cite{dhiman2023samplesize,litjens2017survey}. Cross-validation and related resampling strategies can make more efficient use of limited data, but they do not remove the uncertainty caused by small sample size.

Generalization also depends on whether future data resemble the development data. CT examinations may vary across scanners, acquisition protocols, reconstruction settings, and patient populations \cite{rizzo2018radiomics}. Evaluation in a small cohort should therefore be interpreted as evidence about performance under the available data conditions, not as proof of clinical robustness. The theoretical point is that regression evaluation must be patient-level and that generalization claims require caution under limited data.
